\documentclass{article}

\usepackage[utf8]{inputenc}
\usepackage{amsmath, esint}
\usepackage{wasysym}
\usepackage{qrcode}
\usepackage[colorlinks]{hyperref}
\usepackage{lmodern}
\usepackage{graphicx}
\usepackage{xcolor}
\usepackage[left=2cm, top=3cm, right=2cm]{geometry}
\usepackage{minted}
\usepackage{booktabs}
\usepackage{svg}
\usepackage{xcolor}
\usepackage{tabularx}
\definecolor{LightGray}{gray}{0.975}
\hypersetup{
  colorlinks,
  linkcolor=blue,
  urlcolor=blue
}

\title{ML101 \\ Lab 02: A `gentle' Introduction to Linear Model Selection \& Regularization.}
\author{Andrés Oswaldo Calderón Romero, Ph.D.}
\date{\today}

\begin{document}

\maketitle

\section{Introduction}
\label{sec:intro}

While Ordinary Least Squares (OLS) linear regression offers an intuitive baseline for predictive modeling, its performance degrades when confronting high-dimensional feature spaces, multicollinearity, or scarce observations relative to predictors ($p \gg n$). In such regimes, unconstrained least squares estimates suffer from high variance, causing severe overfitting and poor out-of-sample generalization. To mitigate these shortcomings, modern statistical learning relies on model selection and regularization paradigms: discrete subset selection (best subset, forward, and backward stepwise) isolates parsimonious predictor sets, shrinkage methods like Ridge ($\ell_2$) and Lasso ($\ell_1$) constrain coefficients to trade small amounts of bias for major variance reductions, and dimension reduction approaches like Principal Components Analysis project data onto orthogonal directions of maximal variance.

Building upon the foundations presented in Chapter~6 of \textit{An Introduction to Statistical Learning} (ISLRv2), this lab focuses on implementing, diagnosing, and contrasting these workflows on real-world empirical data. By operationalizing these techniques outside of the native R ecosystem (e.g., in Python), students will explore practical data pitfalls and evaluate the critical trade-offs among predictive accuracy, computational tractability, and model interpretability.

% \section{Additional Resources}
% \label{sec:resources}

\section{Lab: Linear Models and Regularization}
\label{sec:lab}

In this lab, we will work through the exercises in Chapter~6 of ISLRv2\footnote{James, G., Witten, D., Hastie, T., \& Tibshirani, R. (2023). \textit{An Introduction to Statistical Learning: with Applications in R} (2nd ed.). Springer. \url{https://www.statlearning.com/}}. Specifically, from Section~\textit{6.5 Lab: Linear Models and Regularization Methods} (p.~267), we will explore the following subsections:

\begin{itemize}
\item 6.5.1 Subset Selection Methods,
\item 6.5.2 Ridge Regression and the Lasso, and
\item 6.5.3 PCR and PLS Regression.
\end{itemize}

You do not need to include evidence from this walkthrough in your lab report; however, you must familiarize yourself with these tools to complete the autonomous work in the following section.

\section{Autonomous Work}
\label{sec:work}

So, now that you are familiar with model selection and regularization techniques, you will apply them to a dataset of your choice. You may revisit the dataset used in Lab 01, choose a new benchmark from the UCI Machine Learning Repository, or --we highly encouraged-- use your own data. Once you have selected your dataset, apply the methods covered in Section~\ref{sec:lab}. Specifically, you will perform:

\begin{enumerate}
  \item Best Subset Selection.
  \item Forward Stepwise Selection.
  \item Backward Stepwise Selection.
  \item Ridge Regression.
  \item Lasso Regression.
  \item Principal components analysis.
\end{enumerate}

Do not be discouraged if a technique fails or yields poor performance. The primary objective is to familiarize yourself with the available tools and learn how to diagnose and address data-related pitfalls. Please note: you must use a programming language other than R (e.g., Python or Julia).

\section{Deliverables}
Each group must submit a well-structured \textbf{PDF} report documenting the implementation and results of the methods described in Section~\ref{sec:work}. Your report should thoroughly discuss key findings, performance trade-offs, and noteworthy observations. Submissions must be uploaded through the designated Brightspace\texttrademark{} assignment link within \textbf{21 days} of this lab session.

\vspace{5mm}
Happy Hacking! \includesvg[width=4mm]{figures/sunglasses}

\end{document}

